% !TeX program = tectonic
\documentclass[11pt,a4paper]{article}
\usepackage{fontspec}
\usepackage[english]{babel}
\babelfont{rm}[Language=Default]{Liberation Serif}
\babelfont[Language=Default]{sf}{Carlito}
\babelfont[Language=Default]{tt}{DejaVu Sans Mono}
\usepackage{amsmath,amssymb,amsthm}
\usepackage{geometry}
\geometry{a4paper, top=2.4cm, bottom=2.4cm, left=2.4cm, right=2.4cm}
\usepackage{booktabs,tabularx,enumitem,xcolor,tcolorbox}
\tcbuselibrary{breakable,skins}
\definecolor{linkblue}{HTML}{1F4E79}
\definecolor{statgreen}{HTML}{2F6B3A}
\definecolor{goldacc}{HTML}{8B7E5A}
\usepackage[colorlinks=true,linkcolor=linkblue,urlcolor=linkblue,unicode=true]{hyperref}
\theoremstyle{plain}
\newtheorem{theorem}{Theorem}
\newtheorem{lemma}[theorem]{Lemma}
\newtheorem{proposition}[theorem]{Proposition}
\newtheorem{corollary}[theorem]{Corollary}
\theoremstyle{definition}
\newtheorem{definition}[theorem]{Definition}
\theoremstyle{remark}
\newtheorem{remark}[theorem]{Remark}
\newtcolorbox{statusbox}[1][]{colback=statgreen!5,colframe=statgreen!75,
  title={\textbf{Summary of results:} #1},fonttitle=\small,boxrule=0.7pt,arc=2pt,breakable}
\newtcolorbox{databox}[1][]{colback=goldacc!6,colframe=goldacc!80,
  title={\textbf{Protocol data:} #1},fonttitle=\small,boxrule=0.7pt,arc=2pt,breakable}
\newtcolorbox{verifybox}[1][]{colback=linkblue!5,colframe=linkblue!80,
  title={\textbf{Verification:} #1},fonttitle=\small,boxrule=0.7pt,arc=2pt,breakable}
\widowpenalty=10000
\clubpenalty=10000
\setlength{\parskip}{0.35em}
\newcommand{\CC}{\mathbb{C}}
\newcommand{\RR}{\mathbb{R}}
\newcommand{\ZZ}{\mathbb{Z}}
\newcommand{\QQ}{\mathbb{Q}}
\newcommand{\Ga}{\Gamma}
\newcommand{\Bfun}{\mathrm{B}}
\newcommand{\im}{\operatorname{Im}}
\newcommand{\re}{\operatorname{Re}}
\newcommand{\lcm}{\operatorname{lcm}}
\newcommand{\SNF}{\operatorname{SNF}}
\newcommand{\Jac}{\operatorname{Jac}}
\newcommand{\rank}{\operatorname{rank}}
\newcommand{\Ind}{\operatorname{Index}}
\newcommand{\disc}{\operatorname{disc}}
\begin{document}
\begin{center}
{\LARGE\bfseries Certificate H and the Klein Stand}\\[0.4em]
{\large the Gross Gamma-normalization, the volume 1125, and the heptad arithmetic}\\[0.6em]
{\normalsize Isaev Iskhak Khamzatovich}\\[0.2em]
{\normalsize The <<Dynamic Principle>> Program --- the \texttt{hodge-laboratory}}\\[0.15em]
{\small Standalone theorem monograph · Монография 16/16}
\end{center}
\vspace{0.4em}
{\color{linkblue}\hrule height 0.8pt}
\vspace{0.8em}

\section{Setting}

Certificate H --- the summit of the program arch: the full Gross
$\Ga$-normalization on the $N=15/30$ stand, closing the
$\Ga$-remainder of certificate G. The third stand --- the Klein
quartic Jacobian with its arithmetic --- is assembled here as well.

\begin{theorem}[H1 --- the cyclotomic lattice]\label{th:H1stand}
The CM lattice: $\QQ(\zeta_{30})=\QQ(\zeta_{15})$ --- one stand,
two $\Ga$-levels. The Riemann form is integral (Ramanujan sums);
the polarization type (SNF) is $(1,1,5,5,15,15,15,15)$; the
discriminant is $3^4\cdot5^6=1265625=1125^2$.
\end{theorem}

\begin{proof}
The Ramanujan form is integral; the SNF is computed by the exact
integer algorithm; the product of invariant factors
$1\cdot1\cdot5\cdot5\cdot15^4=1265625$ matches $3^4\cdot5^6$ ---
verified by factorization. The field: $\phi(30)=\phi(15)=8$; the
isomorphism $\zeta_{30}\mapsto-\zeta_{15}^8$.
\end{proof}

\begin{theorem}[H2 --- the $\Ga$-volume]\label{th:H2stand}
\[
  \mathrm{vol}_h=(2\pi)^{32}\prod_{k=1}^{N-1}\Ga(k/N)^{-c_k}=1125
  \qquad (N=15\ \text{and}\ 30),
\]
the exponents $c_k$ exact (for $N=15$: $c_k=4$ at ten values of $k$,
$c_k=6$ at $k=3,6,9,12$); the residual $1.01\cdot10^{-119}$;
$\det V^2=1265625/256$; PSLQ deflation ($N=15$) recovers the integer
exponent vector; $N=30$ certified by the sine count.
\end{theorem}

\begin{proof}
The Chowla--Selberg product over characters: the exponents collect
the conductor multiplicities $h_d$. Computation at $dps=120$ gives
$1125$ with residual $10^{-119}$; the agreement with
$\sqrt{\disc}$ of H1 is exact. The sine count:
$\Ga(k/30)\Ga((30-k)/30)=\pi/\sin(\pi k/30)$ closes the identity at
even $k$ without PSLQ.
\end{proof}

\begin{theorem}[H3--H6]\label{th:H36}
\emph{(H3)} Every Fermat period is the $\Ga$-monomial
$\frac1N\zeta^{ra+sb}\Ga(a/N)\Ga(b/N)/\Ga((a+b)/N)$ ($91/406$
periods; quadratures $10^{-121}$); level $30$ with odd $k$ is not
reducible to $15$.
\emph{(H4)} The moduli $|\omega(\tau_i)^2|$ are $\Ga$-monomials; the
phase $\psi=\frac{(1+3\sqrt5)+i(\sqrt{15}-\sqrt3)}8=R/|R|$ is
recovered by LLL; $|R|=3-\sqrt5$ is a unit of $\QQ(\sqrt5)$.
\emph{(H5)} The ladder $\pi/7$, $\pi/15$, $\pi/30$ --- rungs of one
$\mu_N$-quantum ladder; the weights $4\sin(\pi k/N)$ agree.
\emph{(H6)} The denominator bound $Q_{\text{stand}}=2N\max s_i=480$;
the F rationalization closes without a scan.
\end{theorem}

\begin{proof}
(H3) --- the closed period form + certificate B ($69$ tests,
$3.9\cdot10^{-36}$) + deep quadratures ($10^{-121}$);
non-reducibility: $\Ga(1/30)$ is not reduced by the inventory to
denominator $15$. (H4) --- LLL in $\QQ(\sqrt{-3},\sqrt5)$ recovers
the exact coordinates; $R\bar R=(3-\sqrt5)^2$ exact.
(H5) --- the phases $\pi/7$ (Klein), $\pi/15$, $\pi/30$ are levels
of the construction $\delta=\pi/N$; the weights agree with the
Ramanujan tables. (H6) --- $2\cdot30\cdot8=480$; all measurements
inside the F1 uniqueness intervals.
\end{proof}

\begin{databox}[Klein --- the stand within H]
Smoothness: $28(xyz)^3=0$; genus $3$; $\Delta=-343=-7^3$;
$c_4=105$; $j=-3375=-15^3$; the form $W^2=v^3-35v-98$; roots $7,
\frac{-7\pm\sqrt{-7}}2$; $\Omega=1.9333117056\ldots$ (three
schemes, agreement $1.4\cdot10^{-32}$); $\tau=(1+\sqrt{-7})/2$;
$j(\tau)=-3375$ from the period ($57$ digits, anchor $j(i)=1728$);
$\zeta+\zeta^2+\zeta^4=\tau-1$; isotypy $(1,1,0,1)$;
$\lambda_1$(Heawood)$=\sqrt2$, multiplicity $6$; the Fano code
$C_7+\mu_4$.
\end{databox}

\begin{statusbox}[summary]
Certificate H accepted ($6/6$): the cyclotomic lattice, the volume
$1125$ in the full $\Ga$-normalization, the $\Ga$-monomial periods,
the Gross normalization, the ladder, the denominators --- the arch
of the program is closed. Runtime $\sim$2 minutes.
\end{statusbox}

\begin{verifybox}[reproduction]
\texttt{python3 laboratory.py} $\to$ ``Certificates $\to$ H'';
``Stands $\to$ Klein''; Lean: \texttt{stand\_disc},
\texttt{snf\_product}, \texttt{klein\_integers}.
\end{verifybox}

\vfill
{\color{linkblue}\hrule height 0.6pt}
\vspace{0.5em}
{\small\color{gray}
License: individual exclusive license of Isaev Iskhak Khamzatovich.
All rights to the computations, formulas, and texts belong to the author.
Verification: \texttt{python3 laboratory.py} (``Stands''/``Certificates''),
Lean: \texttt{verification/lean/HodgeLaboratory.lean}.
}
\end{document}
